% GENERATED FILE. Edit canonical lessons, then run npm run generate:books.
% canonical-source-hash: fnv1a64:5c04de83d7737ba6
% canonical-lessons: PA-C227-aukha, PA-C227-malna, PA-C227-jhagg, PA-C227-mail, PA-C227-tuti

\chapter{Soap and Water}
\label{ch:pa-a2-soap-and-water}
\hlchaptermodality{227}

\begin{chapteropening}
\textbf{By the end of this chapter:} \emph{I can talk about washing my hands and face.}

Complete the last lesson of chapter 227: I can talk about washing my hands and face.
\end{chapteropening}

\section[\texorpdfstring{\pa{ਔਖਾ} (aukhā)}{ਔਖਾ (aukhā)}]{\pa{ਔਖਾ} (aukhā) — hard to do; difficult}

\label{lesson:PA-C227-aukha}



\begin{quote}
Before the new one: say the Punjabi for to shut a door, then the Punjabi for a walk.
\end{quote}

\subsection*{You'll want to know: \pa{ਔਖਾ}}
\textbf{\pa{ਔਖਾ}} — \emph{aukhā} — \textquotedblleft{}hard to do; difficult\textquotedblright{}.

\begin{quote}
\pa{ਸਰਦੀ} \pa{ਵਿੱਚ} \pa{ਸਵੇਰੇ} \pa{ਨਹਾਉਣਾ} \pa{ਔਖਾ} \pa{ਹੈ।}
\end{quote}

\begin{quote}
\emph{sardī vich savere nahāuṇā aukhā hai.}
\end{quote}

\begin{quote}
Bathing early on a winter morning is hard.
\end{quote}

\begin{cousinweb}
\textbf{\pa{ਔਖਾ}} (\emph{aukhā}) is inherited: CDIAL sets it under Sanskrit \emph{apakṣa}, \textquotedblleft{}without wings, without followers\textquotedblright{}, and lists Punjabi \emph{aukhā}, \textquotedblleft{}difficult, distressed\textquotedblright{} (CDIAL 414).
\end{cousinweb}

\begin{grammarlens}[title={What you've built}]
One more word for saying what you do.
\end{grammarlens}

\subsection*{Your turn}
\begin{itemize}
\raggedright
  \item \emph{Say it:} \emph{aukhā}
  \item \emph{Say it:} \emph{aukhā}, once more
  \item \emph{Say it:} \emph{sardī vich savere nahāuṇā aukhā hai.} — Bathing early on a winter morning is hard.
  \item \emph{Recall:} say the Punjabi for to shut a door, then the Punjabi for a walk, then say \emph{aukhā} again
\end{itemize}

\begin{tcolorbox}[breakable,colback=teal!4,colframe=teal!35!black,title={Before you move on}]
What does \pa{ਔਖਾ} mean? (\textquotedblleft{}Hard to do; difficult\textquotedblright{}.) Say it once more.
\end{tcolorbox}
\section[\texorpdfstring{\pa{ਮਲਣਾ} (malṇā)}{ਮਲਣਾ (malṇā)}]{\pa{ਮਲਣਾ} (malṇā) — to rub}

\label{lesson:PA-C227-malna}



\begin{quote}
Before the new one: say the Punjabi for a walk, then the Punjabi for hard to do.
\end{quote}

\subsection*{You'll want to know: \pa{ਮਲਣਾ}}
\textbf{\pa{ਮਲਣਾ}} — \emph{malṇā} — \textquotedblleft{}to rub\textquotedblright{}.

\begin{quote}
\pa{ਮੈਂ} \pa{ਹੱਥਾਂ} \pa{ਉੱਤੇ} \pa{ਸਾਬਣ} \pa{ਮਲਦਾ} \pa{ਹਾਂ।}
\end{quote}

\begin{quote}
\emph{maiṁ hatthāṁ utte sābaṇ maldā hāṁ.}
\end{quote}

\begin{quote}
I rub soap on my hands.
\end{quote}

Past (I …): \pa{ਮੈਂ} \pa{ਸਾਬਣ} \pa{ਮਲਿਆ}, \emph{maiṁ sābaṇ maliā}.

\begin{cousinweb}
\textbf{\pa{ਮਲਣਾ}} (\emph{malṇā}) is inherited from a reconstructed \emph{*marati}, \textquotedblleft{}crushes, rubs\textquotedblright{}; CDIAL lists Punjabi \emph{maḷnā}, \textquotedblleft{}to rub\textquotedblright{} (CDIAL 9870).
\end{cousinweb}

\begin{grammarlens}[title={What you've built}]
One more word for saying what you do.
\end{grammarlens}

\subsection*{Your turn}
\begin{itemize}
\raggedright
  \item \emph{Say it:} \emph{malṇā}
  \item \emph{Say it:} \emph{malṇā}, once more
  \item \emph{Say it:} \emph{maiṁ hatthāṁ utte sābaṇ maldā hāṁ.} — I rub soap on my hands.
  \item \emph{Recall:} say the Punjabi for a walk, then the Punjabi for hard to do, then say \emph{malṇā} again
  \item \emph{Write it:} copy \textbf{\pa{ਮਲਣਾ}} once, checking the whole word
\end{itemize}

\begin{tcolorbox}[breakable,colback=teal!4,colframe=teal!35!black,title={Before you move on}]
What does \pa{ਮਲਣਾ} mean? (\textquotedblleft{}To rub\textquotedblright{}.) Say it once more.
\end{tcolorbox}
\section[\texorpdfstring{\pa{ਝੱਗ} (jhagg)}{ਝੱਗ (jhagg)}]{\pa{ਝੱਗ} (jhagg) — lather; foam}

\label{lesson:PA-C227-jhagg}



\begin{quote}
Before the new one: say the Punjabi for hard to do, then the Punjabi for to rub.
\end{quote}

\subsection*{You'll want to know: \pa{ਝੱਗ}}
\textbf{\pa{ਝੱਗ}} — \emph{jhagg} — \textquotedblleft{}lather; foam\textquotedblright{}.

\begin{quote}
\pa{ਸਾਬਣ} \pa{ਦੀ} \pa{ਝੱਗ} \pa{ਚਿੱਟੀ} \pa{ਹੈ।}
\end{quote}

\begin{quote}
\emph{sābaṇ dī jhagg chiṭṭī hai.}
\end{quote}

\begin{quote}
The soap lather is white.
\end{quote}

A feminine noun.

\begin{cousinweb}
\textbf{\pa{ਝੱਗ}} (\emph{jhagg}) is inherited from a reconstructed \emph{*jhaggā}, \textquotedblleft{}froth, scum\textquotedblright{}; CDIAL lists Punjabi \emph{jhagg}, \textquotedblleft{}foam, froth, scum\textquotedblright{} (CDIAL 5322).
\end{cousinweb}

\begin{grammarlens}[title={What you've built}]
One more word for saying what you do.
\end{grammarlens}

\subsection*{Your turn}
\begin{itemize}
\raggedright
  \item \emph{Say it:} \emph{jhagg}
  \item \emph{Say it:} \emph{jhagg}, once more
  \item \emph{Say it:} \emph{sābaṇ dī jhagg chiṭṭī hai.} — The soap lather is white.
  \item \emph{Recall:} say the Punjabi for hard to do, then the Punjabi for to rub, then say \emph{jhagg} again
  \item \emph{Write it:} copy \textbf{\pa{ਝੱਗ}} once, checking the whole word
\end{itemize}

\begin{tcolorbox}[breakable,colback=teal!4,colframe=teal!35!black,title={Before you move on}]
What does \pa{ਝੱਗ} mean? (\textquotedblleft{}Lather; foam\textquotedblright{}.) Say it once more.
\end{tcolorbox}
\section[\texorpdfstring{\pa{ਮੈਲ} (mail)}{ਮੈਲ (mail)}]{\pa{ਮੈਲ} (mail) — dirt; grime}

\label{lesson:PA-C227-mail}



\begin{quote}
Before the new one: say the Punjabi for to rub, then the Punjabi for lather.
\end{quote}

\subsection*{You'll want to know: \pa{ਮੈਲ}}
\textbf{\pa{ਮੈਲ}} — \emph{mail} — \textquotedblleft{}dirt; grime\textquotedblright{}.

\begin{quote}
\pa{ਮੇਰੇ} \pa{ਪੈਰਾਂ} \pa{ਉੱਤੇ} \pa{ਮੈਲ} \pa{ਹੈ।}
\end{quote}

\begin{quote}
\emph{mere pairāṁ utte mail hai.}
\end{quote}

\begin{quote}
There is dirt on my feet.
\end{quote}

A feminine noun.

\begin{cousinweb}
\textbf{\pa{ਮੈਲ}} (\emph{mail}) is inherited from a reconstructed \emph{*malin}, \textquotedblleft{}dirty\textquotedblright{}; CDIAL lists Punjabi \emph{mail}, \textquotedblleft{}dirt\textquotedblright{} (CDIAL 9904).
\end{cousinweb}

\begin{grammarlens}[title={What you've built}]
One more word for saying what you do.
\end{grammarlens}

\subsection*{Your turn}
\begin{itemize}
\raggedright
  \item \emph{Say it:} \emph{mail}
  \item \emph{Say it:} \emph{mail}, once more
  \item \emph{Say it:} \emph{mere pairāṁ utte mail hai.} — There is dirt on my feet.
  \item \emph{Recall:} say the Punjabi for to rub, then the Punjabi for lather, then say \emph{mail} again
  \item \emph{Write it:} copy \textbf{\pa{ਮੈਲ}} once, checking the whole word
\end{itemize}

\begin{tcolorbox}[breakable,colback=teal!4,colframe=teal!35!black,title={Before you move on}]
What does \pa{ਮੈਲ} mean? (\textquotedblleft{}Dirt; grime\textquotedblright{}.) Say it once more.
\end{tcolorbox}
\section[\texorpdfstring{\pa{ਟੂਟੀ} (ṭūṭī)}{ਟੂਟੀ (ṭūṭī)}]{\pa{ਟੂਟੀ} (ṭūṭī) — a tap; a spout}

\label{lesson:PA-C227-tuti}



\begin{quote}
Before the new one: say the Punjabi for lather, then the Punjabi for grime.
\end{quote}

\subsection*{You'll want to know: \pa{ਟੂਟੀ}}
\textbf{\pa{ਟੂਟੀ}} — \emph{ṭūṭī} — \textquotedblleft{}a tap; a spout\textquotedblright{}.

\begin{quote}
\pa{ਮੈਂ} \pa{ਟੂਟੀ} \pa{ਖੋਲ੍ਹਦੀ} \pa{ਹਾਂ} \pa{ਅਤੇ} \pa{ਮੂੰਹ} \pa{ਧੋਂਦੀ} \pa{ਹਾਂ।}
\end{quote}

\begin{quote}
\emph{maiṁ ṭūṭī kholhdī hāṁ ate mūnh dhondī hāṁ.}
\end{quote}

\begin{quote}
I turn on the tap and wash my face.
\end{quote}

A feminine noun.

\begin{cousinweb}
\textbf{\pa{ਟੂਟੀ}} (\emph{ṭūṭī}) is inherited from Sanskrit \emph{tuṇḍa}, \textquotedblleft{}beak, snout\textquotedblright{}; CDIAL lists Punjabi \emph{ṭūṭī}, \textquotedblleft{}spout\textquotedblright{} (CDIAL 5853).
\end{cousinweb}

\begin{grammarlens}[title={What you've built}]
One more word for saying what you do.
\end{grammarlens}

\subsection*{Your turn}
\begin{itemize}
\raggedright
  \item \emph{Say it:} \emph{ṭūṭī}
  \item \emph{Say it:} \emph{ṭūṭī}, once more
  \item \emph{Say it:} \emph{maiṁ ṭūṭī kholhdī hāṁ ate mūnh dhondī hāṁ.} — I turn on the tap and wash my face.
  \item \emph{Recall:} say the Punjabi for lather, then the Punjabi for grime, then say \emph{ṭūṭī} again
  \item \emph{Write it:} copy \textbf{\pa{ਟੂਟੀ}} once, checking the whole word
\end{itemize}

\begin{tcolorbox}[breakable,colback=teal!4,colframe=teal!35!black,title={Before you move on}]
What does \pa{ਟੂਟੀ} mean? (\textquotedblleft{}A tap; a spout\textquotedblright{}.) Say it once more.
\end{tcolorbox}
